\subsection*{Appendix C: Term-by-term expansion of the Podolsky operator, with every derivative}

Section~4 quantized the classical Hamiltonian (3.22) by applying the Podolsky operator (4.5) to the ro-vibrational metric (4.15) and rescaling to the flat measure. This appendix supplies what the main text compressed: the expansion of the quantized kernel of (4.41),
\begin{equation}
\hat T_{\mathrm{kernel}}=\tfrac12\sum_{\alpha\beta}\hat A_\alpha\,\mu_{\alpha\beta}(Q)\,\hat A_\beta,\qquad
\hat A_\alpha\equiv\hat J_\alpha-\hat L_\alpha-\hat\pi_\alpha,\qquad
\hat L_\alpha\equiv\hat L_{\mathrm{elec},\alpha}, \tag{C.1}
\end{equation}
into its nine products, with \emph{every} coordinate derivative written out and the Hermiticity of each block verified from first principles. All arguments use only: the adjoint of a product reverses the order, $(\hat X\hat Y)^{\dagger}=\hat Y^{\dagger}\hat X^{\dagger}$; integration by parts on the measure $d\mu_W$; and the product rule. No commutator algebra of body-fixed angular momenta appears anywhere (cf.\ the Remark in \S\,4.11). Repeated Cartesian indices $\alpha,\beta\in\{1,2,3\}$ and mode indices $b,c\in\{1,\dots,f\}$ are summed.

\subsubsection*{C.1\quad Toolkit}

The three momentum-like operators live in three separate coordinate sectors:
\begin{equation}
\hat J_\alpha=-i\hbar\,(\mathsf E^{-1})_{s\alpha}\partial_{\Omega_s}\ \ \text{[eq.~(4.24), Euler angles]},\quad
\hat P_b=-i\hbar\,\partial_{Q_b}\ \ \text{[modes]},\quad
\hat L_\alpha=-i\hbar\,\epsilon_{\alpha\beta\gamma}\textstyle\sum_i\bar r_{i\beta}\partial_{\bar r_{i\gamma}}\ \ \text{[electrons]}. \tag{C.2}
\end{equation}
Because they differentiate disjoint sets of variables, \emph{operators from different sectors commute}, and each sector's functions are constants for the other sectors' derivatives. Each is Hermitian on its factor of $d\mu_W$: $\hat J_\alpha$ on $\sin\theta\,d\Omega$ by the divergence identity (4.26) [via (4.27)], $\hat L_\alpha$ on the flat electron measure by (4.18), and $\hat P_b$ on the flat $\prod_bdQ_b$ trivially. The only non-commutativity in the problem is the product rule within the $Q$-sector,
\begin{equation}
\hat P_b\,f(Q)=f(Q)\,\hat P_b-i\hbar\,(\partial_{Q_b}f)
\qquad\Longrightarrow\qquad
\hat P_b\,\mu_{\alpha\beta}=\mu_{\alpha\beta}\hat P_b-i\hbar\,\partial_b\mu_{\alpha\beta},\quad
\hat P_b\,\sigma_{\alpha c}=\sigma_{\alpha c}\hat P_b-i\hbar\,\zeta^{\alpha}_{bc}, \tag{C.3}
\end{equation}
where the last relation uses $\sigma_{\alpha b}=\sum_a\zeta^{\alpha}_{ab}Q_a$ [(3.17a)], linear in $Q$ with \emph{constant} $\zeta^{\alpha}_{ab}=\bigl(\sum_j\mathbf l_{ja}\times\mathbf l_{jb}\bigr)_\alpha$, so $\partial_{Q_b}\sigma_{\alpha c}=\zeta^{\alpha}_{bc}$. The antisymmetry of the cross product gives the sum rule that recurs throughout:
\begin{equation}
\zeta^{\alpha}_{ab}=-\zeta^{\alpha}_{ba}\qquad\Longrightarrow\qquad
\zeta^{\alpha}_{bb}=0,\qquad \sum_b\partial_{Q_b}\sigma_{\alpha b}=\sum_b\zeta^{\alpha}_{bb}=0 . \tag{C.4}
\end{equation}
\textbf{$\hat\pi_\alpha$ is Hermitian on the flat mode measure.} With $\hat\pi_\alpha=\sum_b\sigma_{\alpha b}\hat P_b$, take the adjoint (order reverses; $\sigma$ real, $\hat P_b$ Hermitian), then slide $\hat P_b$ back to the right with (C.3), then apply (C.4):
\begin{equation}
\hat\pi_\alpha^{\dagger}
=\sum_b\hat P_b\,\sigma_{\alpha b}
=\sum_b\bigl(\sigma_{\alpha b}\hat P_b-i\hbar\,\zeta^{\alpha}_{bb}\bigr)
=\hat\pi_\alpha . \tag{C.5}
\end{equation}
Had $\sum_b\zeta^\alpha_{bb}$ been nonzero, $\hat\pi_\alpha$ would have required explicit symmetrization; the antisymmetry of $\zeta$ spares us. Finally, $\mu_{\alpha\beta}(Q)=\mu_{\beta\alpha}(Q)$ is a real symmetric function of $Q$ alone.

\subsubsection*{C.2\quad The nine products}

Substitute $\hat A_\alpha=\hat J_\alpha-\hat L_\alpha-\hat\pi_\alpha$ into (C.1) and multiply out, \emph{keeping the operator order} (the factors need not commute):
\begin{align}
\hat A_\alpha\mu_{\alpha\beta}\hat A_\beta
=\;&\underbrace{\hat J_\alpha\mu_{\alpha\beta}\hat J_\beta}_{\text{(T1) rotation}}
+\underbrace{\hat L_\alpha\mu_{\alpha\beta}\hat L_\beta}_{\text{(T2) electronic}}
+\underbrace{\hat\pi_\alpha\mu_{\alpha\beta}\hat\pi_\beta}_{\text{(T3) vibrational}}
-\underbrace{\bigl(\hat J_\alpha\mu_{\alpha\beta}\hat L_\beta+\hat L_\alpha\mu_{\alpha\beta}\hat J_\beta\bigr)}_{\text{(T4) rotation--electron}}\nonumber\\
&-\underbrace{\bigl(\hat J_\alpha\mu_{\alpha\beta}\hat\pi_\beta+\hat\pi_\alpha\mu_{\alpha\beta}\hat J_\beta\bigr)}_{\text{(T5) rotation--vibration}}
+\underbrace{\bigl(\hat L_\alpha\mu_{\alpha\beta}\hat\pi_\beta+\hat\pi_\alpha\mu_{\alpha\beta}\hat L_\beta\bigr)}_{\text{(T6) vibration--electron}} . \tag{C.6}
\end{align}
Each cross term collects the two equal classical contributions $X\mu Y$ and $Y\mu X$, which is why it appears as a symmetric sum of orderings. We now treat the six blocks one at a time: for each we verify Hermiticity and write the explicit differential form.

\subsubsection*{C.3\quad Diagonal terms}

\paragraph{(T1)\ \ $\hat J_\alpha\mu_{\alpha\beta}\hat J_\beta$.}
\emph{Hermiticity.} Take the adjoint ($\mu$ real, $\hat J_\alpha^\dagger=\hat J_\alpha$ on $\sin\theta\,d\Omega$; order reverses), then relabel the summed dummy indices $\alpha\leftrightarrow\beta$, then use $\mu_{\beta\alpha}=\mu_{\alpha\beta}$:
\begin{equation}
\bigl(\hat J_\alpha\mu_{\alpha\beta}\hat J_\beta\bigr)^{\dagger}
=\hat J_\beta\,\mu_{\alpha\beta}\,\hat J_\alpha
=\hat J_\alpha\,\mu_{\beta\alpha}\,\hat J_\beta
=\hat J_\alpha\,\mu_{\alpha\beta}\,\hat J_\beta . \tag{C.7}
\end{equation}
Hermitian exactly as written --- no symmetrization, no commutator evaluated. \emph{Differential form.} $\mu(Q)$ commutes with the Euler derivatives, and inserting (C.2),
\begin{equation}
\hat J_\alpha\mu_{\alpha\beta}\hat J_\beta
=-\hbar^{2}\mu_{\alpha\beta}\,(\mathsf E^{-1})_{s\alpha}\Bigl[(\mathsf E^{-1})_{t\beta}\,\partial_{\Omega_s}\partial_{\Omega_t}
+\bigl(\partial_{\Omega_s}(\mathsf E^{-1})_{t\beta}\bigr)\,\partial_{\Omega_t}\Bigr], \tag{C.8}
\end{equation}
with the trigonometric entries of $\mathsf E^{-1}$ read off (4.12). Every term carries at least one Euler derivative: the rotational block generates \emph{no} pure multiplicative term, in agreement with the Laplace--Beltrami computation (4.29).

\paragraph{(T2)\ \ $\hat L_\alpha\mu_{\alpha\beta}\hat L_\beta$.}
Word-for-word the same as (T1) with $\hat J\to\hat L$: the adjoint chain (C.7) uses only $\hat L_\alpha^\dagger=\hat L_\alpha$ [(4.18)] and the symmetry of $\mu$. Hermitian as written; all terms carry electron derivatives; no multiplicative residue.

\paragraph{(T3)\ \ $\hat\pi_\alpha\mu_{\alpha\beta}\hat\pi_\beta$.}
\emph{Hermiticity.} Use $\hat\pi^{\dagger}=\hat\pi$ from (C.5), then relabel $\alpha\leftrightarrow\beta$ and use $\mu$ symmetric --- the same three-step chain as (C.7). Hermitian as written, but \emph{only} because $\hat\pi$ is Hermitian, i.e.\ ultimately because of the antisymmetry (C.4). \emph{Differential form.} Write $\hat\pi_\alpha=\sigma_{\alpha b}\hat P_b$ and slide the left momentum past the $Q$-function $\mu_{\alpha\beta}\sigma_{\beta c}$ with (C.3):
\begin{align}
\hat\pi_\alpha\mu_{\alpha\beta}\hat\pi_\beta
&=\sigma_{\alpha b}\,\hat P_b\,\bigl(\mu_{\alpha\beta}\sigma_{\beta c}\bigr)\,\hat P_c
=\sigma_{\alpha b}\Bigl[\bigl(\mu_{\alpha\beta}\sigma_{\beta c}\bigr)\hat P_b-i\hbar\,\partial_b\!\bigl(\mu_{\alpha\beta}\sigma_{\beta c}\bigr)\Bigr]\hat P_c\nonumber\\[2pt]
&=\underbrace{G_{bc}\,\hat P_b\hat P_c}_{\text{2nd order}}
\;-\;i\hbar\,\underbrace{\sigma_{\alpha b}\Bigl[(\partial_b\mu_{\alpha\beta})\sigma_{\beta c}+\mu_{\alpha\beta}\,\zeta^{\beta}_{bc}\Bigr]\hat P_c}_{\text{1st order}},
\qquad G_{bc}\equiv(\boldsymbol\sigma^{\mathsf T}\boldsymbol\mu\boldsymbol\sigma)_{bc}, \tag{C.9}
\end{align}
the Coriolis-dressed vibrational operator. Every surviving term carries at least one $\hat P$: again no pure multiplicative term. \emph{Equivalence with the ordered form.} The same sum rule shows that (T3) equals the fully left-ordered operator: moving the left $\hat P_b$ past its own coefficient $\sigma_{\alpha b}$,
\begin{equation}
\hat P_b\,G_{bc}\,\hat P_c
-\hat\pi_\alpha\mu_{\alpha\beta}\hat\pi_\beta
=\bigl[\hat P_b,\sigma_{\alpha b}\bigr]\,\mu_{\alpha\beta}\sigma_{\beta c}\hat P_c
=-i\hbar\Bigl(\textstyle\sum_b\zeta^{\alpha}_{bb}\Bigr)\mu_{\alpha\beta}\sigma_{\beta c}\hat P_c=0 , \tag{C.10}
\end{equation}
so $\tfrac12\hat\pi_\alpha\mu_{\alpha\beta}\hat\pi_\beta=\tfrac12\hat P_b\,G_{bc}\,\hat P_c$ \emph{exactly}. This identity is what lets the sandwich-ordered kernel of (4.41) match the Laplace--Beltrami vibrational block (4.31), whose dressed part is precisely $\tfrac12\hat P_bG_{bc}\hat P_c$ after the rescaling (see \S\,C.5).

\subsubsection*{C.4\quad Cross terms}

\paragraph{(T4)\ \ $\hat J_\alpha\mu_{\alpha\beta}\hat L_\beta+\hat L_\alpha\mu_{\alpha\beta}\hat J_\beta$.}
All three factors act on disjoint variables and mutually commute. A single ordering is already Hermitian,
\begin{equation}
\bigl(\hat J_\alpha\mu_{\alpha\beta}\hat L_\beta\bigr)^{\dagger}
=\hat L_\beta\,\mu_{\alpha\beta}\,\hat J_\alpha
=\mu_{\alpha\beta}\,\hat J_\alpha\hat L_\beta
=\hat J_\alpha\mu_{\alpha\beta}\hat L_\beta , \tag{C.11}
\end{equation}
and the two orderings in the sum are equal (relabel $\alpha\leftrightarrow\beta$ and commute), so
\begin{equation}
\hat J_\alpha\mu_{\alpha\beta}\hat L_\beta+\hat L_\alpha\mu_{\alpha\beta}\hat J_\beta
=2\,\mu_{\alpha\beta}\,\hat J_\alpha\hat L_\beta
\qquad(\text{Hermitian; no residue}). \tag{C.12}
\end{equation}
This is the rotation--electron Coriolis coupling; no $\hat P$ appears, so there is no differential subtlety.

\paragraph{(T5)\ \ $\hat J_\alpha\mu_{\alpha\beta}\hat\pi_\beta+\hat\pi_\alpha\mu_{\alpha\beta}\hat J_\beta$.}
Now the subtlety: $\hat\pi$ does \emph{not} commute with $\mu$, since by (C.3) $\hat\pi_\beta\mu_{\alpha\beta}=\mu_{\alpha\beta}\hat\pi_\beta-i\hbar\,\sigma_{\beta b}\partial_b\mu_{\alpha\beta}\neq\mu_{\alpha\beta}\hat\pi_\beta$. A single ordering is therefore \emph{not} Hermitian by itself; its adjoint is exactly the \emph{other} ordering ($\hat J$ commutes with $\mu$ and $\hat\pi$; relabel $\alpha\leftrightarrow\beta$):
\begin{equation}
\bigl(\hat J_\alpha\mu_{\alpha\beta}\hat\pi_\beta\bigr)^{\dagger}
=\hat\pi_\beta\,\mu_{\alpha\beta}\,\hat J_\alpha
=\hat\pi_\alpha\,\mu_{\alpha\beta}\,\hat J_\beta . \tag{C.13}
\end{equation}
The cross term is of the form $\hat X+\hat X^{\dagger}$, automatically Hermitian --- \emph{this is why the classical product $-2J\mu\pi$ must be quantized as the symmetric sum}, and it is precisely the symmetric sum that the Laplace--Beltrami operator produced by itself in (4.30)/(4.35). Pulling $\hat J$ out front and expanding the anticommutator with (C.3),
\begin{equation}
\hat J_\alpha\mu_{\alpha\beta}\hat\pi_\beta+\hat\pi_\alpha\mu_{\alpha\beta}\hat J_\beta
=\hat J_\alpha\bigl(\mu_{\alpha\beta}\hat\pi_\beta+\hat\pi_\beta\mu_{\alpha\beta}\bigr)
=\hat J_\alpha\Bigl(\underbrace{2\,\mu_{\alpha\beta}\sigma_{\beta c}\hat P_c}_{\text{1st order in }\hat P}
\;\underbrace{-\,i\hbar\,\sigma_{\beta c}\,\partial_c\mu_{\alpha\beta}}_{\text{0th order in }\hat P}\Bigr). \tag{C.14}
\end{equation}
The first piece is the rotation--vibration Coriolis coupling ($\hat J$ times $\hat P$); the second is the imaginary, $\hat J$-linear counterterm that makes the sum Hermitian. It is \emph{not} a stand-alone multiplicative residue --- it still carries $\hat J_\alpha$ --- consistent with the Laplace--Beltrami result (4.35) that the Coriolis block contributes nothing to the pseudopotential. As a cross-check of (C.14) against the pre-rescaling form (4.30): expanding the second group of (4.30) by the product rule,
\begin{equation}
\tfrac{1}{\sqrt\gamma}\,\hat P_b\,\sqrt\gamma\;\sigma_{\alpha b}\mu_{\alpha\beta}\,\hat J_\beta
=\Bigl[\sigma_{\alpha b}\mu_{\alpha\beta}\hat P_b
-i\hbar\,\partial_b\!\bigl(\sigma_{\alpha b}\mu_{\alpha\beta}\bigr)
-\tfrac{i\hbar}{2}\,(\partial_b\ln\gamma)\,\sigma_{\alpha b}\mu_{\alpha\beta}\Bigr]\hat J_\beta , \tag{C.15}
\end{equation}
and conjugating by $\gamma^{1/4}$ as in \S\,4.9 cancels the $\partial_b\ln\gamma$ terms between the two groups, reproducing exactly $-\tfrac12\hat J_\alpha\{\mu_{\alpha\beta},\hat\pi_\beta\}$ of (4.35), i.e.\ $-\tfrac12\times$(C.14).

\paragraph{(T6)\ \ $\hat L_\alpha\mu_{\alpha\beta}\hat\pi_\beta+\hat\pi_\alpha\mu_{\alpha\beta}\hat L_\beta$.}
Identical to (T5) with $\hat J\to\hat L$ ($\hat L$ also commutes with $\mu$ and $\hat\pi$):
\begin{equation}
\hat L_\alpha\mu_{\alpha\beta}\hat\pi_\beta+\hat\pi_\alpha\mu_{\alpha\beta}\hat L_\beta
=\hat L_\alpha\Bigl(2\,\mu_{\alpha\beta}\sigma_{\beta c}\hat P_c-i\hbar\,\sigma_{\beta c}\,\partial_c\mu_{\alpha\beta}\Bigr), \tag{C.16}
\end{equation}
the vibration--electron angular coupling.

\subsubsection*{C.5\quad The measure residue $U_{\mathrm{Cor}}$, every derivative}

Section 4.9 derived the identity-block residue $U_{\mathrm{vib}}$ [(4.34)] in full and deferred the residue of the dressed block. Here it is. Conjugate the dressed part of the vibrational Laplace--Beltrami block (4.31) by $\gamma^{1/4}$ and act on $\psi_W$; each line is one product rule [$G_{bc}=(\boldsymbol\sigma^{\mathsf T}\boldsymbol\mu\boldsymbol\sigma)_{bc}$, symmetric]:
\begin{align}
\partial_{Q_c}\!\bigl(\gamma^{-1/4}\psi_W\bigr)&=\gamma^{-1/4}\Bigl[\partial_{Q_c}\psi_W-\tfrac14(\partial_c\ln\gamma)\psi_W\Bigr],\nonumber\\
\gamma^{1/2}\,G_{bc}\,\partial_{Q_c}\!\bigl(\gamma^{-1/4}\psi_W\bigr)&=\gamma^{1/4}\,G_{bc}\Bigl[\partial_{Q_c}\psi_W-\tfrac14(\partial_c\ln\gamma)\psi_W\Bigr],\nonumber\\
\partial_{Q_b}\Bigl(\gamma^{1/4}\,G_{bc}\bigl[\cdots\bigr]\Bigr)
&=\gamma^{1/4}\Bigl[G_{bc}\,\partial_b\partial_c\psi_W+(\partial_bG_{bc})\,\partial_c\psi_W
+\tfrac14(\partial_b\ln\gamma)G_{bc}\,\partial_c\psi_W-\tfrac14 G_{bc}(\partial_c\ln\gamma)\,\partial_b\psi_W\nonumber\\
&\qquad\quad-\tfrac14\,\partial_b\bigl(G_{bc}\,\partial_c\ln\gamma\bigr)\psi_W
-\tfrac{1}{16}\,G_{bc}(\partial_b\ln\gamma)(\partial_c\ln\gamma)\,\psi_W\Bigr]. \tag{C.17}
\end{align}
The two first-order $\ln\gamma$ terms cancel by the symmetry of $G_{bc}$ (relabel $b\leftrightarrow c$), leaving $G\partial\partial+(\partial G)\partial$ --- which is exactly $-\tfrac{2}{\hbar^2}\times\tfrac12\hat P_bG_{bc}\hat P_c$ --- plus a pure multiplicative remainder. Multiplying by $-\tfrac{\hbar^2}{2}$ (the prefactors $\gamma^{1/4}\cdot\gamma^{-1/2}\cdot\gamma^{1/4}=1$ combine to unity),
\begin{equation}
\gamma^{1/4}\,\hat T_{\mathrm{dress}}\,\gamma^{-1/4}
=\tfrac12\,\hat P_b\,G_{bc}\,\hat P_c+U_{\mathrm{Cor}},\qquad
\boxed{\;U_{\mathrm{Cor}}=\frac{\hbar^{2}}{8}\Bigl[\partial_b\bigl(G_{bc}\,\partial_c\ln\gamma\bigr)+\tfrac14\,G_{bc}\,(\partial_b\ln\gamma)(\partial_c\ln\gamma)\Bigr].\;} \tag{C.18}
\end{equation}
By (C.10), $\tfrac12\hat P_bG_{bc}\hat P_c=\tfrac12\hat\pi_\alpha\mu_{\alpha\beta}\hat\pi_\beta$, so the operator part is exactly the (T3) sandwich of the kernel (4.41), and the entire ordering debt of the dressed block is the multiplicative $U_{\mathrm{Cor}}$. Together with $U_{\mathrm{vib}}$ of (4.34) [the same computation with $G_{bc}\to\delta_{bc}$], the total residue is
\begin{equation}
\hat U=U_{\mathrm{vib}}+U_{\mathrm{Cor}}
=\frac{\hbar^{2}}{8}\Bigl[\partial_b\bigl(g^{QQ}_{bc}\,\partial_c\ln\gamma\bigr)+\tfrac14\,g^{QQ}_{bc}\,(\partial_b\ln\gamma)(\partial_c\ln\gamma)\Bigr],
\qquad g^{QQ}_{bc}=\delta_{bc}+G_{bc}, \tag{C.19}
\end{equation}
every ingredient of which reduces, via Jacobi's formula (4.36) and the linearity (4.37), to the constant $\boldsymbol\zeta$'s and the linear $\boldsymbol\sigma(Q)$; the Watson sum rules then collapse (C.19) to the trace $\hat U=-\tfrac{\hbar^2}{8}\sum_\alpha\mu_{\alpha\alpha}$ of (4.38)~\cite{watson1968}.

\textbf{The division of labour is clean.} In the explicit forms (C.8), (C.9), (C.12), (C.14), (C.16), every surviving kernel term carries at least one derivative operator ($\hat P$, $\hat J$, or $\hat L$): \emph{the ordering of the kernel never generates a pure multiplicative term}. The pseudopotential (4.38) is an $\mathcal O(\hbar^2)$ $c$-number that originates entirely in the \emph{curved integration measure} $\sqrt\gamma$, extracted by the rescaling $\psi_W=\gamma^{1/4}\psi$ --- (C.18) is its dressed-block share. Kernel ordering $\rightarrow$ the Hermitian operators of this appendix; the measure $\rightarrow$ the pseudopotential.

\subsubsection*{C.6\quad Worked example: a bent triatomic}

To make ``every derivative'' concrete, take a bent triatomic $XY_2$ (e.g.\ H${}_2$O), a $C_{2v}$ asymmetric top with $N_{\mathrm{nucl}}=3$ and $f=3N_{\mathrm{nucl}}-6=3$ normal coordinates: the symmetric stretch $Q_1$ and bend $Q_2$ (both $a_1$) and the antisymmetric stretch $Q_3$ ($b_2$). Choose the body frame as the instantaneous principal-axis system, with $a,b$ in the molecular plane and $c\perp$ to it.

\paragraph{Inertia and its derivatives.}
At equilibrium the inertia tensor is diagonal, $\mathbf I_N^{\mathrm{eq}}=\operatorname{diag}(I_a^{e},I_b^{e},I_c^{e})$, and because the equilibrium mass distribution is planar the perpendicular-axis theorem gives the exact relation
\begin{equation}
I_c^{e}=I_a^{e}+I_b^{e}. \tag{C.20}
\end{equation}
The totally symmetric coordinates $Q_1,Q_2$ preserve the $C_{2v}$ symmetry, so to linear order they only \emph{scale} the diagonal moments, whereas the antisymmetric $Q_3$ tilts the principal axes and enters $\mathbf I_N$ only at second order. Hence
\begin{equation}
I_\alpha(Q)=I_\alpha^{e}+a_\alpha^{1}Q_1+a_\alpha^{2}Q_2+\mathcal O(Q^2),
\qquad a_\alpha^{b}\equiv\frac{\partial I_\alpha}{\partial Q_b}\Big|_{\mathrm{eq}}, \tag{C.21}
\end{equation}
with the inertial derivatives $a_\alpha^{b}$ computable from the mode vectors $\mathbf l_{jb}$ evaluated at equilibrium. These are exactly the derivatives that populate the rotation--vibration Jacobian $\sqrt{\gamma(Q)}$ of (4.21).

\paragraph{Coriolis coupling.}
The only nonzero constant Coriolis matrix in $C_{2v}$ couples the two $a_1$ modes to the $b_2$ mode about the $b$ axis (the axis whose rotation maps a symmetric into an antisymmetric displacement):
\begin{equation}
\boldsymbol\sigma_1=\zeta^{b}_{13}\,Q_3\,\hat{\mathbf b},\quad
\boldsymbol\sigma_2=\zeta^{b}_{23}\,Q_3\,\hat{\mathbf b},\quad
\boldsymbol\sigma_3=-\bigl(\zeta^{b}_{13}Q_1+\zeta^{b}_{23}Q_2\bigr)\hat{\mathbf b},
\qquad \zeta^{b}_{ab}=-\zeta^{b}_{ba}, \tag{C.22}
\end{equation}
so $\boldsymbol\sigma\boldsymbol\sigma^{T}=\bigl[(\zeta^{b}_{13}Q_3)^2+(\zeta^{b}_{23}Q_3)^2+(\zeta^{b}_{13}Q_1+\zeta^{b}_{23}Q_2)^2\bigr]\hat{\mathbf b}\hat{\mathbf b}^{T}$ corrects only the $bb$ element of the inertia,
\begin{equation}
\mathbf I'(Q)=\operatorname{diag}\!\bigl(I_a(Q),\,I_b(Q)-s(Q),\,I_c(Q)\bigr),\qquad
s(Q)\equiv(\zeta^{b}_{13}Q_3)^2+(\zeta^{b}_{23}Q_3)^2+(\zeta^{b}_{13}Q_1+\zeta^{b}_{23}Q_2)^2 . \tag{C.23}
\end{equation}
Note $s(Q)=\mathcal O(Q^2)$, so it does not affect the linear-order derivatives (C.21) but does enter the pseudopotential and the exact Jacobian.

\paragraph{Determinant, volume element, and its explicit derivatives.}
With $\mathbf I'$ diagonal in (C.23),
\begin{equation}
\gamma(Q)=\det\mathbf I'=I_a(Q)\,\bigl(I_b(Q)-s(Q)\bigr)\,I_c(Q),\qquad
\sqrt{g}=\sin\theta\,\sqrt{I_a\,(I_b-s)\,I_c}, \tag{C.24}
\end{equation}
and the vibrational Jacobian carries the explicit derivatives
\begin{equation}
\partial_{Q_b}\ln\gamma
=\frac{a_a^{b}}{I_a}+\frac{a_b^{b}-\partial_{Q_b}s}{I_b-s}+\frac{a_c^{b}}{I_c},
\qquad
\partial_{Q_3}s=2\bigl[(\zeta^{b}_{13})^2+(\zeta^{b}_{23})^2\bigr]Q_3,\;\;
\partial_{Q_1}s=2\zeta^{b}_{13}\bigl(\zeta^{b}_{13}Q_1+\zeta^{b}_{23}Q_2\bigr),\ \text{etc.}, \tag{C.25}
\end{equation}
which are precisely the first-order remainders exhibited under (4.31) and removed by the $\gamma^{1/4}$ rescaling of \S\,4.9.

\paragraph{Pseudopotential.}
Since $\mathbf I'$ is diagonal, $\boldsymbol\mu=\operatorname{diag}\bigl(1/I_a,\,1/(I_b-s),\,1/I_c\bigr)$ and the trace formula (4.38) is fully explicit:
\begin{equation}
\boxed{\;\hat U(Q)=-\frac{\hbar^2}{8}\Bigl[\frac{1}{I_a(Q)}+\frac{1}{I_b(Q)-s(Q)}+\frac{1}{I_c(Q)}\Bigr].\;}\tag{C.26}
\end{equation}
Expanding to the order at which it first varies with the coordinates, and using (C.20),
\begin{equation}
\hat U(Q)=-\frac{\hbar^2}{8}\Bigl(\frac{1}{I_a^{e}}+\frac{1}{I_b^{e}}+\frac{1}{I_c^{e}}\Bigr)
+\frac{\hbar^2}{8}\sum_{b=1,2}\Bigl(\frac{a_a^{b}}{(I_a^{e})^2}+\frac{a_b^{b}}{(I_b^{e})^2}+\frac{a_c^{b}}{(I_c^{e})^2}\Bigr)Q_b+\mathcal O(Q^2). \tag{C.27}
\end{equation}
The constant is an inconsequential shift of the electronic--vibrational zero; the linear term is a genuine $\hbar^2$ force on the normal coordinates --- the physical content of the Watson pseudopotential for this molecule --- produced entirely by differentiating $\sqrt{\gamma(Q)}$ through (C.25). Every step from the metric (4.15) to the surface (C.27) used only $\partial_\phi,\partial_\theta,\partial_\chi$ and $\partial_{Q_b}$.

\subsubsection*{C.7\quad Summary}

\begin{enumerate}
\item The quantized kernel (C.1) expands into the nine products (C.6). Blocks built from mutually commuting, individually Hermitian factors --- (T1) $\hat J\mu\hat J$, (T2) $\hat L\mu\hat L$, (T4) $\hat J\mu\hat L$ --- are Hermitian in a single ordering, by the adjoint-reversal argument (C.7) alone.
\item Blocks pairing $\hat\pi$ with $\hat J$ or $\hat L$ --- (T5), (T6) --- are Hermitian only as the symmetric sum $\hat X+\hat X^\dagger$, which is exactly the ordering the Laplace--Beltrami operator produces by itself [(4.30), (4.35)]. (T3) $\hat\pi\mu\hat\pi$ is the fortunate middle case, Hermitian because $\hat\pi$ is [(C.5)], and equal to the left-ordered $\tfrac12\hat P G\hat P$ by the sum rule (C.10).
\item No kernel ordering produces a pure multiplicative term. The Watson pseudopotential comes solely from the measure: $U_{\mathrm{vib}}$ [(4.34)] plus $U_{\mathrm{Cor}}$ [(C.18)] collapse to $-\tfrac{\hbar^2}{8}\sum_\alpha\mu_{\alpha\alpha}$ [(4.38)].
\item Every Hermiticity statement above used only integration by parts, the product rule, and the divergence identities (4.18), (4.26), (C.4) --- never a commutator of body-fixed angular momenta.
\end{enumerate}
